\SourceLeaf{308}
\fancyhead[L]{\small INVESTIGATIONS INTO TOPICS IN HIGHER GEODESY}
\renewcommand{\Article}[1]{\par\medskip\Needspace{4\baselineskip}\pdfbookmark[1]{Art. #1}{geodaesie2-art-#1}\centerline{\textbf{#1.}}\nobreak\smallskip}
\begin{center}
INVESTIGATIONS\par\bigskip
INTO\par\bigskip
{\large TOPICS IN HIGHER GEODESY.}\par\medskip
\rule{.65in}{.3pt}
\end{center}
\noindent\hspace*{\parindent}Given the length of a triangle side on the Earth's surface and the azimuth and geographical latitude at one endpoint, finding the azimuth and latitude at the other endpoint and the difference in longitude between the two points is one of the main tasks of higher geodesy. For the sphere, the relation between these six quantities was given at the end of the first memoir in the simplest form suitable for calculation to the highest accuracy. That form also lends itself readily to a convenient solution of the problem itself. This naturally makes one want a similar method of solution that applies directly to the ellipsoid, and the purpose of the present memoir is to develop one. First, however, the solution for the sphere will be set out more clearly. To make references easier, I continue the numbering of the articles from the first memoir.

\Article{18}
To judge more accurately how much precision the formulas in Art. 17 give, we must work out the terms of the next order. It is worth
